\section{The topological Tate calculation for the subgroup of order three}
\label{sec:subgroup}

We now determine the part of the $H$-homotopy fixed point spectral sequence
needed for ascent to $G$, where
\[
 H=\langle\sigma\rangle=C_3\subset G=\langle g\rangle=C_9,
 \qquad \sigma=g^3.
\]
Hill's report on joint work with Hopkins and Ravenel
\cite[pp.~2699--2700, Propositions~2--3]{Hill2007} sketches the
$C_p$ calculation at heights $(p-1)f$, conditional on Hopkins's
coefficient-action conjecture. Here the required finite-page statements
follow from the completed coefficient model and fixed geometry of
Section~\ref{sec:coefficients}, together with the finite Euler bound of
Section~\ref{sec:detector}.

Throughout this section, $E_r(H)$ denotes the \emph{topological}
Euler-localized homotopy fixed point spectral sequence. Its differentials have
bidegree $(r,r-1)$ in cohomological filtration and internal degree. The
complete-local coefficient spectral sequence of
Section~\ref{sec:coefficients} has already been evaluated; none of its
differentials is being relabelled as a differential of $E_r(H)$.

We recall the notation and consequences of
Theorems~\ref{coef:model}, \ref{coef:fixed-geometry}, and
\ref{coef:hom-vanishing}. We have
\begin{equation}
 \label{h:e2}
 E_2(H)=K[q^{\pm1},b^{\pm1}]\otimes_K\Lambda_K(V),
 \qquad K=k[[n_1,n_2]],\qquad k=\F_{3^6},
\end{equation}
where
\[
 |q|=(0,-6),\qquad |b|=(2,0),\qquad |V|=(1,4).
\]
The free $K$-module $V$ of rank three has a $T=\F_{27}^{\times}$-weight
basis $v,\theta,y$ of weights $-2,0,6$, respectively. The class $\theta$ is
the restriction of the first strict coordinate-change cocycle, with a
nonzero scalar normalization permitted. In particular, it is fixed by the
centralizer $\mathcal C=C_{\mathcal O_D^\times}(H)$. Modulo
$\mathfrak n=(n_1,n_2)$, the action of $g$ on $V$ is one Jordan block:
$\theta$ is fixed, $g(y)-y$ is a nonzero multiple of $\theta$, and
$g(v)-v$ has a nonzero $y$-coefficient. The tangent weights of $K$ are $6$
and $-2$.

The coefficient classes in $K$ and the class
\begin{equation}
 \label{h:a}
 A=qb^{-2},\qquad |A|=(-4,-6),
\end{equation}
are cycles on every page on which they occur, by
Lemma~\ref{coef:scalar-cycles}. Here $A$ is a unit on the nonzero
localized pages. This assertion concerns the norm class divided by an
Euler square; it does not assert that $q$ or $b$ is separately a permanent
cycle.

Let $h_1$ be the first Hasse section on the fixed formal-module
deformation. As in Section~\ref{sec:coefficients}, put
\begin{equation}
 \label{h:curve-data}
 J=h_1^3\in Kq^{-2},\qquad F=q^2J\in K,\qquad
 S_0=K/(F),\qquad W=(V/K\theta)\otimes_K S_0.
\end{equation}
The element $F$ is a regular parameter. Write $P_n$ for the scalar line
$Kq^n$, or its base change to $S_0$ when working on the curve.
On the reduced height-two curve, whose ring is denoted by $D'$,
the second Hasse section $h$ has Lie line $w^{-8}$ and is a parameter
after trivializing that line. The inclusion $S_0\simeq(D')^3$ and the
evaluation of $W$ give the following rules:
\begin{enumerate}[label=(\roman*)]
 \item There is no $\mathcal C$-equivariant $S_0$-linear homomorphism $W\to P_n$,
 even after passing to meromorphic sections.
 \item An invariant integral vector in $WP_n$ evaluates, if nonzero, to
 a constant multiple of $h^d$, where
 \begin{equation}
 \label{h:hasse-rule}
  8d=2-3n,\qquad d\geq0.
 \end{equation}
 \item If $h$ occurs primitively in $WP_{-2}$, it defines a stable line
 $\ell=hP_2\subset W$. After extension to $D'$, the quotient
 $W/\ell$ has character $w^{-15}\chi$, where $\chi$ is a nontrivial
 sign character. A meromorphic invariant section with this sign twist
 does not exist.
\end{enumerate}
The two evaluated basis values of $W$, after the common Lie factor is
removed, have orders $0$ and $1$. Their coefficients lie in the cube
ring, so their possible leading orders have different residues modulo
three. We will use this fact when dividing an invariant vector by a
power of the curve parameter.

Only differential lengths congruent to $1$ modulo $4$ occur. Indeed,
the central stabilizer element $-1$ acts trivially on the ungraded
deformation ring and by $-1$ on its Lie line. It therefore acts on
$E_t$ by $(-1)^{t/2}$. Evenness first forces $r$ odd, and naturality
for this central action then forces $(-1)^{(r-1)/2}=1$ for a nonzero
$d_r$. Thus $r\equiv1\pmod4$. Moreover,
Proposition~\ref{det:euler-bound} implies that the unit of $E_r(H)$ is
hit by a differential of length at most $157$. Indeed, the actual Euler
class has $79$th power zero, its power symbol has filtration $158$, and
the finite skeletal tower then gives an incoming differential of length
at most $158$; the congruence on lengths improves this to $157$. All
uses of this bound below are statements about finite pages.

\subsection{The first two topological differentials}

\begin{lemma}
 \label{h:greek-evaluation}
 With compatible nonzero scalar normalizations, the images of the
 indicated Greek-letter classes are
 \[
  \alpha_1=\theta,\qquad
  \beta_1=bJ,\qquad
  \beta_{3/3}=bJ^3,\qquad
  \beta_4=\pm bJ^5.
 \]
 In particular,
 \begin{equation}
 \label{h:d5-b}
  d_5(b)=\theta b^3,\qquad d_5(\theta)=0.
 \end{equation}
\end{lemma}

\begin{proof}
Evaluate the integral Bockstein representatives on the fixed deformation
over $D_0=W(k)[\zeta_3][[r_1,r_2]]$. The coefficient lines in the three
even degrees under consideration have trivial $H$-action. In the cyclic
resolution, the connecting map for reduction modulo three is therefore
$b$ times the value of the intermediate one-cocycle at $\sigma$.
This is an integral connecting-map calculation before reduction by the
uniformizer $\zeta_3-1$.

Modulo that uniformizer, the formal-module relation
$(\sigma-1)^2=-3\sigma$ gives
\[
 v_1=0,\qquad v_2=-t_1(\sigma)^4.
\]
The intermediate $v_1$-Bocksteins for $v_2$, $v_2^3$, and $v_2^4$
therefore have values, respectively,
\[
 t_1(\sigma)^3,\qquad t_1(\sigma)^9,\qquad
 \pm t_1(\sigma)^{15}.
\]
For example, the last value is the reduction of
$4v_2^3t_1(\sigma)^3$, and the middle value follows by cubing the
right-unit difference before dividing by $v_1^3$. These calculations
can use Hazewinkel generators; changing to the Araki parameters of the
fixed-geometry calculation leaves these reductions unchanged.
The scalar-cohomology evaluation in
Theorem~\ref{coef:fixed-geometry} is injective, giving the asserted
classes in \eqref{h:e2}.

We use the classical Adams--Novikov differentials
\begin{equation}
 \label{h:classical-differentials}
 d_5(\beta_{3/3})=\text{unit}\cdot\alpha_1\beta_1^3,
 \qquad
 d_9(\alpha_1\beta_4)=\text{unit}\cdot\beta_1^6,
\end{equation}
and the permanence of $\alpha_1$ and $\beta_1$; see
\cite[Lemma 2.5 and Table 1]{BelmontShimomura}.
The comparison of Section~\ref{sec:detector} carries these formulas to
the present spectral sequence. Since
$\beta_{3/3}=b^{-2}\beta_1^3$, the first formula gives
$d_5(b)=\theta b^3$ after normalizing $\theta$. Cancellation by
$\beta_1^3$ is valid on this page: its scalar factor is nonzero in the
domain $K[q^{\pm1},b^{\pm1}]$, and the exterior algebra is free over
that domain. Permanence of $\alpha_1$ gives $d_5(\theta)=0$.
\end{proof}

We first determine $d_5$ on the remaining exterior classes without
assuming its value. Write
\[
 d_5(z)=b^2\Phi(z),\qquad
 \Phi:V\longrightarrow\Lambda^2V.
\]
The map is $K$-linear because $K$ consists of cycles. Modulo
$\mathfrak n$, the $T$-weights allow only
$v\mapsto\lambda_v\theta v$ and
$y\mapsto\lambda_y\theta y$. Equivariance for the residual Jordan
action of $g$ gives $\lambda_v=\lambda_y$. There are no linear
coefficient corrections: comparing the weights of $v,y$ with those of
$\theta v,\theta y,vy$ shows that none of the required additional
weights is $6$ or $-2$. Thus
\begin{equation}
 \label{h:phi-initial}
 \Phi\equiv\lambda_0\,\theta\wedge(-)\pmod{\mathfrak n^2},
 \qquad \lambda_0\in k.
\end{equation}

Set
\[
 \widehat V=b^{-1}V,\qquad \Theta=b^{-1}\theta,
 \qquad \mu=\lambda_0-1.
\]
On $K[A^{\pm1},b^{\pm1}]\otimes\Lambda(\widehat V)$, the normalized
derivation $d=b^{-3}d_5$ satisfies
\begin{equation}
 \label{h:normalized-d5}
 db=b\Theta,\qquad d\Theta=0,\qquad
 d|_{\widehat V}=\Psi,\qquad
 \Psi\equiv\mu\Theta\wedge(-)\pmod{\mathfrak n^2}.
\end{equation}
Here $\Psi$ extends as an exterior derivation, and $\Psi^2=0$.
In a fixed $b$-exponent sector $i$, the exterior complex is therefore
\begin{equation}
 \label{h:sector-complex}
 \bigl(\Lambda(\widehat V),\;\Psi+i\Theta\wedge(-)\bigr).
\end{equation}

\begin{lemma}
 \label{h:normalized-d9}
 On the cohomology of \eqref{h:sector-complex}, the operation
 $\kappa=A^2d_9$ preserves both the $A$- and $b$-exponents and lowers
 exterior degree by one. It is a derivation and satisfies
 \begin{equation}
 \label{h:kappa-theta}
  \kappa(\Theta)=cF,\qquad c\in k^{\times}.
 \end{equation}
\end{lemma}

\begin{proof}
The bidegree of $\kappa$ is $(1,-4)$. A monomial
$A^a b^i\Lambda^j\widehat V$ has bidegree
$(-4a+2i-j,-6a+4j)$. If $\Delta a,\Delta i,\Delta j$ are the changes
under $\kappa$, then
\[
 -4\Delta a+2\Delta i-\Delta j=1,
 \qquad -6\Delta a+4\Delta j=-4.
\]
The first equation makes $\Delta j$ odd; the second makes
$\Delta j\equiv-1\pmod3$. Since $0\leq j\leq3$, the only
possibility is $\Delta j=-1$, and then
$\Delta a=\Delta i=0$. Multiplication by the cycle $A^2$ preserves
the derivation rule.

The scalar cohomology of $d_5$ is
$K[A^{\pm1},B^{\pm1}]$, where $B=b^3$, and it has no incoming
$d_5$ boundaries. In particular,
\[
 \beta_1=A^{-2}B^{-1}F
\]
is a non-zero-divisor on this scalar cohomology. The image of
$\alpha_1\beta_4$ is, up to scalar,
$B^{-1}\Theta\beta_1^5$. The second formula of
\eqref{h:classical-differentials}, and the absence of a $d_9$ target
for $B$, therefore give
\[
 B^{-1}d_9(\Theta)\beta_1^5
   =\text{unit}\cdot\beta_1^6.
\]
Cancelling the regular scalar factor proves
$d_9(\Theta)=cA^{-2}F$.
\end{proof}

\begin{lemma}
 \label{h:mu-and-matrix}
 One has $\mu\in\F_3^{\times}$; in integer exponents we use its
 representative $1$ or $-1$. Moreover,
 \[
 \Psi=\mu\Theta\wedge(-)
 \]
 on the complementary two-plane, and
 \begin{align}
 \label{h:e6}
 E_6(H)&=K[A^{\pm1},B^{\pm1}]
       \otimes\Lambda\bigl(\Theta,b^{-1-\mu}(V/K\theta)\bigr),\\
 \label{h:e10}
 E_{10}(H)&=S_0[A^{\pm1},B^{\pm1}]
       \otimes\Lambda\bigl(b^{-1-\mu}W\bigr).
 \end{align}
 The $d_9$ is the Koszul differential on $\Theta,F$ and is zero on the
 other displayed generators.
\end{lemma}

\begin{proof}
Suppose first that $\mu=0$, and put
$z=\Theta\hat v\hat y$ in exponent sector zero. Every boundary in
this sector lies in $\mathfrak n^2$. Since $\Theta z=0$, the
derivation rule and \eqref{h:kappa-theta} give
\begin{equation}
 \label{h:mu-zero-identity}
 \Theta\kappa(z)=cFz
\end{equation}
in cohomology. Choose a $T$-weight representative for $\kappa(z)$.
Its weight is $10$: $z$ has weight $4$ and multiplication by $A^2$
adds weight $6$. The exterior-two basis weights are
\[
 \begin{array}{c|ccc}
  \text{basis vector}&\Theta\hat v&\Theta\hat y&\hat v\hat y\\
  \text{weight}&-2&6&4.
 \end{array}
\]
There is no constant term of weight $10$, and the only possible linear
term is a multiple of $F\hat v\hat y$, because the coefficient tangent
weights are $6,-2$. Equation~\eqref{h:mu-zero-identity} fixes its
coefficient, so
\[
 \kappa(z)\equiv cF\hat v\hat y\pmod{\mathfrak n^2}.
\]
Both $gA/A$ and the determinant of $g$ on $\widehat V$ are $1$ modulo
$\mathfrak n$, and $gF\equiv F\pmod{\mathfrak n^2}$.
Equivariance of $d_9$ therefore makes this representative $g$-fixed
modulo $\mathfrak n^2$. On the other hand, the residual Jordan action
makes $g(\hat v\hat y)-\hat v\hat y$ have a nonzero
$\Theta\hat v$-coefficient. Since $F$ is a parameter, this is a
contradiction. Thus $\mu\ne0$.

Choose a complement to $K\Theta$ and write
\[
 \Psi(z')=\Theta M_0(z')+c_0(z')\hat v\hat y
\]
on that complement, with $M_0\equiv\mu I\pmod{\mathfrak n^2}$.
The exterior-three component of $\Psi^2(z')=0$ is
\[
 c_0\bigl((\operatorname{tr}M_0-M_0)z'\bigr)=0.
\]
The matrix $\operatorname{tr}M_0-M_0$ has invertible reduction
$\mu I$, so $c_0=0$.

If $\mu\notin\F_3$, the complementary exterior-one and
exterior-two terms in every sector are paired with their
$\Theta$-multiples by invertible matrices: their reductions are
$i+\mu$ and $i+2\mu$, for $i\in\F_3$. Only the scalar and
$\Theta$ terms in sector zero remain. The differential
\eqref{h:kappa-theta} then leaves only the even algebra
$S_0[A^{\pm1},B^{\pm1}]$. It has no odd class capable of hitting the
unit on a later page, contradicting the Euler bound. Hence
$\mu\in\F_3^{\times}$.

It remains to remove possible higher coefficient corrections to $M_0$.
In sector $i=-\mu$ put $N=M_0-\mu I$. The complementary
exterior-one and exterior-two cohomology groups are
$\ker N$ and $\coker N$, respectively; multiplication by $\Theta$
is the natural map $\ker N\to\coker N$. For a class
$\alpha\in\coker N$, one has $\Theta\alpha=0$. Applying
$\kappa$ shows that $cF\alpha$ lies in the image of this natural
map. In particular, the two standard cokernel generators have lifts
$k_1,k_2\in\ker N$ such that
\[
 k_j\equiv cF e_j\pmod{\im N}.
\]
Every entry of $N$ lies in $\mathfrak n^2$, so the matrix with columns
$k_1,k_2$ is $cF I$ modulo $\mathfrak n^2$. Its determinant has
nonzero initial form $c^2F^2$. The two columns thus have full rank
over $\operatorname{Frac}K$, and $N=0$.

The cohomology of \eqref{h:sector-complex} now gives
\eqref{h:e6}. The quotient notation $V/K\theta$ is independent of
the chosen complement: changing a lift by a multiple of $\Theta$
changes the corresponding cycle by a boundary in its nonzero sector.
The operation $\kappa$ vanishes on the two complementary generators
because it preserves their nonzero exponent sector, where no scalar
cohomology occurs. It also vanishes on $A$ and $B$. Finally, $F$ is
regular, so taking the Koszul cohomology of
$d_9(\Theta)=cA^{-2}F$ proves \eqref{h:e10}.
\end{proof}

\subsection{Degree restrictions on the remaining plane}

For a component of exterior degree $j\in\{0,1,2\}$, introduce raw
coordinates
\begin{equation}
 \label{h:raw-coordinates}
 m=\frac{s-j}{2},\qquad n=\frac{4j-t}{6}.
\end{equation}
Thus its underlying form is $b^mP_n\Lambda^jW$. It occurs on
\eqref{h:e10} precisely when
\begin{equation}
 \label{h:p-coordinate}
 3p=m+2n+(1+\mu)j,\qquad p\in\Z.
\end{equation}
Here $p$ is the exponent of $B$ after extracting the indicated $A$
and exterior factors. It is not a prime or a Kervaire index.

An increasing differential must have $\Delta j=1$, and a decreasing
one must have $\Delta j=-1$. Solving the two degree equations gives
\begin{equation}
 \label{h:degree-rules}
 \begin{array}{c|ccc}
  r&\Delta j&\Delta n&\Delta p\\\hline
  12l+5&1&-2l&(2l+3+\mu)/3\\
  12l+9&-1&-2l-2&(2l-\mu)/3.
 \end{array}
\end{equation}
The integrality condition in the first row is
$2l+\mu\equiv0\pmod3$, and in the second it is
$2l-\mu\equiv0\pmod3$. These exhaust the possibilities: an odd
filtration change makes $\Delta j$ odd, and the exterior degree is
between zero and two. In the interval $9<r\leq157$ the raw lists are
\begin{equation}
 \label{h:length-table}
 \begin{array}{c|l|l}
  \mu&\text{increasing}&\text{decreasing}\\\hline
  1&17,53,89,125&33,69,105,141\\
  -1&29,65,101,137&21,57,93,129.
 \end{array}
\end{equation}

Before the first further change, decreasing maps are excluded by the
no-homomorphism rule for $W$. This also excludes a map from the
exterior square to the plane, since
$\Hom(\det W,WP_n)\simeq W^{\vee}P_n$.
An increasing map from scalars to the plane requires an invariant
vector in $WP_{-2l}$. The same condition governs maps from the plane
to its determinant, by wedge pairing. Equation~\eqref{h:hasse-rule}
therefore requires
\begin{equation}
 \label{h:first-hasse-degree}
 8d=2+6l,\qquad l=1+4k_0,\qquad d=1+3k_0.
\end{equation}
Combining this with \eqref{h:length-table} and the Euler bound shows
that the first nonzero change is $d_{17}$ for $\mu=1$, or $d_{65}$
for $\mu=-1$.

In either case a primitive vector representing $h$ exists in $WP_{-2}$,
by Theorem~\ref{coef:hom-vanishing}.
Indeed, the nonzero differential supplies a vector evaluating to
$h^{1+3k_0}$. After multiplication by the appropriate power of the
unit $q$, express it in the two evaluation basis values of orders zero
and one. The coefficients are cube series. Since the two possible
leading orders are distinct modulo three, a value of order
$1+3k_0$ forces both coefficients to be divisible by the $k_0$th
power of the scalar parameter cube. Dividing gives a vector evaluating
to $h$, and its order one makes it primitive.

Consequently define
\begin{equation}
 \label{h:xyz}
 \ell=hP_2\subset W,\qquad Z=q^8h^3\in S_0,
 \qquad Y=b^{-1-\mu}q^2h,\qquad X=b^{-1-\mu}v.
\end{equation}
The scalar $Z$ is a parameter of $S_0$. In this notation $h$ denotes
the invariant vector just constructed, whereas $h^3$ denotes the
scalar Hasse section in $P_{-8}$; these are compatible evaluations,
not an exterior cube of a vector.

Invariant meromorphic functions on the curve are constant. Therefore
the first nonzero differential, with the common factor
$A^{-2l}B^{\Delta p}$ suppressed, has the form
\begin{equation}
 \label{h:first-plane-differential}
 \frac{dB}{B}=\epsilon Z^{k_0}Y,
 \qquad dX=\eta Z^{k_0}YX,
 \qquad dY=0,
 \qquad \epsilon,\eta\in k,
\end{equation}
with $\epsilon,\eta$ not both zero. The exterior-square formula is
the wedge pairing with the same invariant line. In exponent $p$, the
scalar-to-$Y$ and $X$-to-$YX$ coefficients are thus
\begin{equation}
 \label{h:first-plane-coefficients}
 p\epsilon Z^{k_0},\qquad(p\epsilon+\eta)Z^{k_0}.
\end{equation}

\begin{lemma}
 \label{h:exclude-minus}
 The case $\mu=-1$ is impossible.
\end{lemma}

\begin{proof}
In this case the first change has length $65$ and $k_0=1$. The only
later possible lengths up to the bound are $93,101,129,137$.
The coefficients in \eqref{h:first-plane-coefficients} are constants
times $Z$. Hence every $E_{66}(H)$ component is a direct sum of free
$S_0$-modules and modules killed by $Z$, and the scalar components
are free. Generically the plane component in a sector is the whole
plane $W$, the line $\ell$, the quotient $W/\ell$, or zero.

At length $93$, the plane-to-scalar map is zero. For the whole plane
this follows from the no-homomorphism rule; for the quotient it
follows from the sign obstruction. For the line, $l=7$ in
\eqref{h:degree-rules}, so a contraction would require a meromorphic
invariant in
\[
 \Hom(P_2,P_{-16})=P_{-18}.
\]
Its Lie exponent is $-54$, which is not divisible by eight. Thus the
map is zero generically, and it is zero integrally because its scalar
target is torsion-free.

Only a decreasing differential can hit the unit, so it would have to
be $d_{129}$. Its source has exterior degree one and $n=22$.
We may choose a $T$-invariant source representative, since $|T|=26$
is invertible in $k$. The two coefficient weights before quotienting
are $3n-2$ and $3n+6$, while $Z$ has weight $-2$. Consequently the
orders of invariant coefficients at $n=22$ satisfy
\[
 66-2-2a\equiv0\pmod{26}
 \quad\hbox{or}\quad
 66+6-2a\equiv0\pmod{26}.
\]
They are respectively $6$ and $10$ modulo $13$. Thus every such
representative on $E_{66}(H)$ is divisible by $Z^6$ within its
component. This statement also holds for subquotient components:
averaging lifts an invariant class to an invariant representative,
and the displayed orders give the divisibility before passage to
the quotient.

There is no outgoing $d_{93}$ on a plane component. Therefore the
$E_{94}(H)$ plane is a quotient of the preceding plane, and every
potential unit source still has a representative $Z^6e$. There may
be incoming $d_{93}$ maps from exterior degree two; their images do
not affect this conclusion. The exterior-two group, in contrast,
is a submodule of its $E_{66}(H)$ predecessor, since no exterior-three
group can map into it. Its torsion is consequently still killed by
$Z$.

Suppose that $Z^6e$ is a $d_{101}$-cycle. Then
$Z^6d_{101}(e)=0$, and the preceding torsion statement gives
$Zd_{101}(e)=0$. Hence $Ze$ is also a cycle, and on $E_{102}(H)$
\[
 [Z^6e]=Z^5[Ze].
\]
There is no change between this page and $E_{129}(H)$. The scalar
groups have received no $d_{93}$ and have only been replaced by
submodules under $d_{101}$; their unit is not divisible by $Z$.
Since $Z$ is a cycle, a $Z$-divisible source cannot map to that unit.
Thus $d_{129}$ cannot hit it. The increasing length $137$ cannot do
so either, contradicting the bound $157$.
\end{proof}

\subsection{The forced contractions and the last differential}

We henceforth have $\mu=1$, so the first change is at length $17$,
$k_0=0$, and $|X|=|Y|=(-3,4)$.

\begin{lemma}
 \label{h:d17-d33}
 One has $\epsilon\ne0$ and $\eta/\epsilon=2$. After rescaling $h$,
 the $d_{17}$ coefficients in \eqref{h:first-plane-coefficients} are
 $p$ and $p-1$. The next nonzero change is $d_{33}$, which contracts
 the surviving $\ell$-generator by a unit times $Z$. Its cohomology
 is
 \begin{equation}
 \label{h:e34}
 E_{34}(H)=k[A^{\pm1},C^{\pm1}]\otimes\Lambda(U),
 \qquad C=B^3=b^9,\qquad U=BX.
 \end{equation}
\end{lemma}

\begin{proof}
All nonzero coefficients of $d_{17}$ in
\eqref{h:first-plane-coefficients} are units, so its cohomology is
torsion-free over $S_0$. If $\epsilon=0$, then $\eta\ne0$ and the
$X,YX$ terms disappear, leaving the exterior algebra on $Y$ over
$S_0[A^{\pm1},B^{\pm1}]$. If $\epsilon\ne0$ but
$\eta/\epsilon\notin\F_3$, the $X,YX$ terms again disappear;
the surviving algebra is exterior on $B^2Y$ over
$S_0[A^{\pm1},B^{\pm3}]$.

In the remaining cases let $p_x\in\{0,1,2\}$ represent
$-\eta/\epsilon\in\F_3$ and set
$C=B^3$. The cohomology is
\begin{equation}
 \label{h:e18-algebra}
 S_0[A^{\pm1},C^{\pm1}]
   \otimes\Lambda(B^2Y,B^{p_x}X).
\end{equation}
Its scalar, line, quotient, and product supports in the $p$ coordinate
are respectively
\begin{equation}
 \label{h:e18-support}
 0,\qquad2,\qquad p_x,\qquad p_x+2\pmod3.
\end{equation}
If $p_x=2$, the two plane columns align and together carry the
corresponding shifted $W$-action. If $p_x\ne2$, the second column
carries the quotient $W/\ell$: the $\ell$ component of a changed
lift lies in a sector killed by the preceding differential.

We check all further operations before any other change. An increasing
map involving the line is the invariant-vector problem for
$\ell P_{-2l}$; after evaluation it has the same condition
$l\equiv1\pmod4$ as \eqref{h:first-hasse-degree}. This also governs
maps from the quotient to the determinant. Increasing maps to the
quotient, or from the line to the determinant, have the nontrivial
sign twist and vanish. The full-plane case has the original
invariant-vector condition. None of the remaining increasing lengths
$53,89,125$ satisfies it.

For a decreasing map, contraction on $\ell=hP_2$ has coefficient
line
\begin{equation}
 \label{h:line-contraction}
 \Hom(P_2,P_{-2l-2})=P_{-2l-4}.
\end{equation}
The Hasse rule makes $l\equiv2\pmod4$; among the decreasing lengths
in \eqref{h:length-table}, only $33$ meets this condition. The same
condition governs contraction from the determinant to the quotient.
The other two line/quotient possibilities have the sign obstruction,
and the aligned full-plane possibilities vanish by the no-homomorphism
rule. Thus the aligned case $p_x=2$ admits no further differential
within the bound and is impossible.

The $d_{33}$ must therefore occur. Its coefficient is a nonzero
constant times $h^3$, by \eqref{h:line-contraction} with $l=2$.
In the algebra \eqref{h:e18-algebra}, this says
\begin{equation}
 \label{h:d33-explicit}
 d_{33}(B^2Y)=aA^{-6}CZ,\qquad
 d_{33}(B^{p_x}X)=0,\qquad a\in k^{\times},
\end{equation}
for $p_x=0$ or $1$. The map on the product follows by the derivation
rule. The factors $A^{-6}$ and $C$ are surviving units. In the
$\epsilon=0$ case the corresponding formula is
$d_{33}(Y)=aA^{-6}BZ$; in the other one-generator case it is the
first formula of \eqref{h:d33-explicit}. Either one-generator case
would consequently leave only an even scalar algebra, incapable of
killing the unit later. This excludes both of those cases and proves
$\epsilon\ne0$ and $p_x\in\{0,1\}$.

Taking cohomology of \eqref{h:d33-explicit} gives
\[
 k[A^{\pm1},C^{\pm1}]\otimes\Lambda(B^{p_x}X).
\]
The remaining odd coefficient has residual $T$-weight $-2$ in the
$W$ line. An increasing map now requires
\[
 3n\equiv3(n-2l)-2\pmod{26},
 \quad\text{equivalently}\quad 6l+2\equiv0\pmod{26},
\]
and a decreasing map requires
\[
 3n-2\equiv3(n-2l-2)\pmod{26},
 \quad\text{equivalently}\quad 6l+4\equiv0\pmod{26}.
\]
Within the bound, these leave only lengths $53$ and $105$,
respectively. Only the latter can hit the unit. At length $105$ the
shift is $\Delta p=5$, so a source in the odd support $p_x$ can
land in scalar support zero only if $p_x=1$. Thus
$\eta/\epsilon=2$, and setting $U=BX$ gives \eqref{h:e34}.
Rescale $h$ to absorb $\epsilon$ in
\eqref{h:first-plane-differential}; then the two $d_{17}$
coefficients become $p$ and $p+2=p-1$.
\end{proof}

\begin{proposition}
 \label{h:unit-death}
 The only possible differential between $E_{34}(H)$ and $E_{105}(H)$
 is
 \begin{equation}
 \label{h:d53}
 d_{53}(C)=e_{53}A^{-8}C^2U,\qquad d_{53}(U)=0,
 \qquad e_{53}\in k.
 \end{equation}
 The class $A^{18}C^{-2}U$ survives to $E_{105}(H)$, and
 \begin{equation}
 \label{h:d105}
 d_{105}(A^{18}C^{-2}U)\in k^{\times}.
 \end{equation}
 In particular, $E_{106}(H)=0$.
\end{proposition}

\begin{proof}
The possible lengths were determined in the proof of
Lemma~\ref{h:d17-d33}. Since
\[
 |C|=(18,0),\qquad |U|=(3,4),\qquad |A|=(-4,-6),
\]
the increasing length $53$, together with the derivation rule, has
exactly the form \eqref{h:d53}. Its scalar is an element of $k$, and
there is no exterior-two group to receive $d_{53}(U)$.

If $e_{53}=0$, the algebra \eqref{h:e34} persists until length
$105$. If $e_{53}\ne0$, then
\[
 d_{53}(C^a)=a e_{53}A^{-8}C^{a+1}U.
\]
The surviving scalar exponents are $a\equiv0\pmod3$, and the
surviving odd exponents are $a\equiv1\pmod3$. Thus
\begin{equation}
 \label{h:e54-nonzero}
 E_{54}(H)=k[A^{\pm1},C^{\pm3}]\otimes\Lambda(CU)
 \qquad(e_{53}\ne0).
\end{equation}
In particular, the notation $C^{-2}U$ still denotes a well-defined
surviving class: in \eqref{h:e54-nonzero} it means
$(C^3)^{-1}(CU)$. Only the surviving Laurent unit $C^3$ is inverted
in this expression.

The bidegree of $A^{18}C^{-2}U$ is $(-105,-104)$, exactly the
bidegree of a possible $d_{105}$ source for the unit. In either case
its source group is one-dimensional. There is no other length after
$33$ and at most $157$ that can hit the unit. The Euler bound
therefore forces \eqref{h:d105}. Since every nonzero scalar is a unit,
the derivation rule makes the cohomology after this differential zero.
Equivalently, if $e_{53}=0$ it pairs all odd and even Laurent slots,
while if $e_{53}\ne0$ it pairs the odd and even slots of
\eqref{h:e54-nonzero}.
\end{proof}

\subsection{The form used in ascent to the group of order nine}

We collect the calculation in raw coordinates so that restriction and
transfer do not require a new choice of generators. On the initial
page allow $j=0,1,2,3$ and set
\begin{equation}
 \label{h:w-coordinate}
 w_0=m+2n+2j.
\end{equation}
After $d_5$, the complementary slots have $w_0=3p$ and their
$\Theta$-multiples have $w_0=3p+1$. After $d_9$, only the former
remain, modulo $F$.

\begin{theorem}
 \label{h:differentials}
 The $H$-Tate calculation has the following properties.
 \begin{enumerate}[label=(\arabic*)]
  \item The first differential is
  \[
  d_5(b^m q^n\xi)
    =(m+2n+2j)b^{m+2}q^n\theta\wedge\xi,
    \qquad \xi\in\Lambda^jV,
  \]
  up to a single nonzero scalar normalization. Its cohomology has
  generators $A,B=b^3,\Theta=b^{-1}\theta$ and
  $b^{-2}(V/K\theta)$ as in \eqref{h:e6}, with $\mu=1$.
  For a fixed $c\in k^{\times}$, the next differential is
  \[
   d_9(\Theta)=cA^{-2}F=cb^4J.
  \]
  Equivalently, on surviving $\Theta$-containing representatives,
  \[
   d_9\bigl([b^m q^n\theta\wedge\xi]\bigr)
     =c[b^{m+5}q^nJ\xi],
  \]
  where $\xi$ is a complementary exterior monomial. These raw
  products are interpreted in the $A,B$ normal form: $A^{-2}$ is
  a surviving unit, while $b^4$ and $J$ need not survive separately.
  Its cohomology is \eqref{h:e10}.
  \item The $d_{17}$ is wedge multiplication by the invariant Hasse
  vector $h$, with common unit shifts and coefficients $p$ on exterior
  degree zero and $p-1$ on exterior degree one. Its survivors have
  supports
  \[
  \begin{array}{c|c}
  \text{component}&p\pmod3\\\hline
  \text{scalars}&0\\
  \ell&2\\
  W/\ell&1\\
  \det W&0.
  \end{array}
  \]
  \item The nonzero $d_{33}$ contracts $\ell$ to scalars and
  $\det W$ to $W/\ell$. Under the identifications by $h$ and
  $h\wedge(-)$, respectively, these maps are multiplication by
  $h^3$, up to nonzero scalar units and the prescribed degree shifts.
  Only scalars at $p\equiv0$ and quotient classes at $p\equiv1$
  remain, now over $S_0/(Z)=k$.
  \item The remaining algebra and differentials are
  \eqref{h:e34}, \eqref{h:d53}, and \eqref{h:d105}. The coefficient
  $e_{53}$ has not been asserted nonzero in this subgroup calculation.
 \end{enumerate}
\end{theorem}

\begin{proof}
The preceding lemmas show that the exact normalized $d_5$ has
$\mu=1$. Thus $d_5(z)=2b^2\theta\wedge z$ for $z\in V$.
Since $A=qb^{-2}$ is a cycle,
$d_5(q)=2qb^2\theta$. The derivation rule gives the coefficient
$m+2n+2j$ in part (1). In raw coordinates, the contraction
$d_9(\Theta)=cA^{-2}F$ is multiplication by $J=q^{-2}F$ after
restoring the corresponding $b$-shifts.

Parts (2) and (3) are Lemma~\ref{h:d17-d33}, with the supports
read from \eqref{h:e18-support} at $p_x=1$. For the product-to-quotient
map, apply the derivation rule to the nonzero line contraction; it
becomes multiplication by the same scalar Hasse cube after using
$h\wedge(-)$ to identify the quotient with the determinant. Part (4)
is Proposition~\ref{h:unit-death}.
\end{proof}

Two consequences specify the interface with the next section.
First, for a Kervaire index $\nu\geq2$, a scalar class in the detector
bidegree has $m=1$ and $n=-2\cdot3^\nu$. Its $d_5$ coefficient is
$1-4\cdot3^\nu\equiv1\pmod3$. A nonzero such scalar cannot be the
restriction of a permanent detector. Second, the unit source in
\eqref{h:d105} has raw coordinates
\begin{equation}
 \label{h:unit-source-raw}
 (j,m,n,w_0)=(1,-53,18,-15).
\end{equation}
The index $n$ is divisible by three. These coordinates will locate
its transfer in the $G$ calculation, where the extra integral
restriction kernels must be retained.
